\subsection*{Problema 44}

\textbf{Enunciado:} Calcular la integral indefinida
\begin{align*}
I = \int x^3 (1 + x^2)^{-1/2} \, dx
\end{align*}

\textbf{Solución:} 
Siguiendo las indicaciones, aplicamos la sustitución $u = 1 + x^2$. 
Derivando ambos lados, obtenemos $du = 2x \, dx$, lo que implica que:
\begin{align*}
x \, dx = \dfrac{1}{2} \, du
\end{align*}

Además, de la definición de $u$, despejamos $x^2$:
\begin{align*}
x^2 = u - 1
\end{align*}

Reescribimos el integrando original separando un factor $x$ para aplicar la sustitución:
\begin{align*}
I &= \int x^2 (1 + x^2)^{-1/2} (x \, dx) \\
I &= \int (u - 1) u^{-1/2} \left(\dfrac{1}{2} \, du\right) \\
I &= \dfrac{1}{2} \int (u - 1) u^{-1/2} \, du
\end{align*}

Multiplicando el término $u^{-1/2}$ por cada componente del binomio, descomponemos la integral en dos potencias simples:
\begin{align*}
I &= \dfrac{1}{2} \int \left( u^{1/2} - u^{-1/2} \right) du
\end{align*}

Integramos cada término utilizando la regla de la potencia:
\begin{align*}
I &= \dfrac{1}{2} \left[ \dfrac{u^{3/2}}{3/2} - \dfrac{u^{1/2}}{1/2} \right] + C \\
I &= \dfrac{1}{2} \left[ \dfrac{2}{3} u^{3/2} - 2 u^{1/2} \right] + C \\
I &= \dfrac{1}{3} u^{3/2} - u^{1/2} + C
\end{align*}

Factorizando $u^{1/2}$ y devolviendo la variable original $u = 1 + x^2$:
\begin{align*}
I &= u^{1/2} \left( \dfrac{1}{3}u - 1 \right) + C \\
I &= (1 + x^2)^{1/2} \left[ \dfrac{1}{3}(1 + x^2) - 1 \right] + C \\
I &= (1 + x^2)^{1/2} \left( \dfrac{x^2 - 2}{3} \right) + C
\end{align*}

Por lo tanto, el resultado de la integral es:
\begin{align*}
I = \dfrac{(x^2 - 2)\sqrt{1 + x^2}}{3} + C
\end{align*}

\begin{figure}[H]
	\centering
	\includegraphics[width=\columnwidth]{semana_05/graficas/SEM05_P44_grafica.png}
	\caption{Representación gráfica de $f(x)$ y su antiderivada $F(x)$ evaluada para la constante de integración $C = 0$ (Problema 44).}
	\label{fig:problema44}
\end{figure}
